Chương này trình bày bối cảnh nghiên cứu, lý do lựa chọn đề tài và định hướng xây dựng hệ thống kiểm soát chất lượng không khí trong căn hộ chung cư. Trọng tâm của khóa luận là kết hợp giám sát môi trường trong nhà, chăm sóc cây trầu bà và điều khiển thiết bị gia dụng để tạo ra một hệ thống IoT có khả năng vận hành thực tế.

\section{Đặt vấn đề}
Chất lượng không khí trong nhà (Indoor Air Quality - IAQ) là một yếu tố có ảnh hưởng trực tiếp đến sức khỏe, khả năng tập trung và sự thoải mái của con người. Theo các khuyến nghị về sức khỏe môi trường, ô nhiễm không khí và điều kiện thông gió kém có thể làm gia tăng nguy cơ khó chịu đường hô hấp, mệt mỏi và suy giảm chất lượng sinh hoạt \cite{who_air_pollution,epa_iaq}. Đối với căn hộ chung cư, không gian sinh hoạt thường tương đối kín, phụ thuộc nhiều vào điều hòa, quạt thông gió và nguồn sáng nhân tạo. Khi mật độ người trong phòng tăng hoặc cửa phòng đóng trong thời gian dài, nồng độ CO2 có thể tăng nhanh, độ ẩm mất cân bằng và nhiệt độ không còn phù hợp. Những thay đổi này không phải lúc nào cũng được người dùng nhận biết kịp thời bằng cảm quan.

Cây trồng trong nhà được sử dụng ngày càng phổ biến vì có giá trị trang trí, tạo cảm giác gần gũi với thiên nhiên và hỗ trợ cải thiện vi khí hậu trong không gian sống. Trong nhóm cây nội thất, cây trầu bà (\textit{Epipremnum aureum}) phù hợp với bối cảnh căn hộ vì dễ nhân giống, sinh trưởng tốt trong ánh sáng lọc, chịu được dải điều kiện môi trường tương đối rộng và được ghi nhận có khả năng hỗ trợ loại bỏ một số chất ô nhiễm trong nhà như formaldehyde, xylene và benzene \cite{pothos_wisconsin}. Các nghiên cứu gần đây về bộ lọc thực vật chủ động cũng sử dụng trầu bà trong cấu hình tường thực vật hoặc hệ thủy canh để xử lý một số VOC như acetone, toluene và $\alpha$-pinene \cite{montaluisa_pothos_biofilter}. Tuy nhiên, cây trồng chỉ phát huy hiệu quả khi được chăm sóc trong điều kiện phù hợp về nước, độ pH, lượng chất hòa tan, ánh sáng, nhiệt độ và độ ẩm. Nếu thiếu hệ thống giám sát, người dùng khó đánh giá đồng thời cả trạng thái cây và chất lượng không khí xung quanh.

Từ thực tế đó, khóa luận đề xuất hệ thống kiểm soát chất lượng không khí trong căn hộ chung cư sử dụng cây trầu bà theo hướng IoT. Hệ thống không chỉ đo đạc các thông số môi trường mà còn liên kết dữ liệu IAQ với điều kiện chăm sóc cây, từ đó hỗ trợ điều khiển quạt, đèn và điều hòa. Kiến trúc gồm một node cảm biến dựa trên Arduino Pro Mini, một gateway ESP32, hai module LoRa để truyền dữ liệu giữa node và gateway, nền tảng ThingsBoard để lưu trữ và hiển thị dữ liệu, cùng ứng dụng Flutter để người dùng theo dõi và điều khiển thiết bị. Ngoài chức năng giám sát, hệ thống có thể điều khiển relay để bật/tắt quạt, đèn và phát lệnh hồng ngoại IR nhằm tăng giảm nhiệt độ điều hòa.
